%!TEX encoding = UTF8
%!TEX root = 

%%%%%%%%%%%%%%%%%%%%%%%%
%%%%%%%%%% SET 1  %%%%%%%%%%
%%%%%%%%%%%%%%%%%%%%%%%% 


\exe{}{
	Calculer la valeur de $6\choose3$ et interpréter le résultat.
}{exe:6choose3}{
	Par définition,
		\[ {6 \choose 3} = \frac{6!}{3!3!} = \frac{6\times5\times4\times3\times2\times1}{(3\times2\times1)^2} = 20. \]
	Interprétation : il existe 20 sous-ensembles de $\{ 1 ; 2 ; 3 ; 4 ; 5 ; 6\}$ de 3 éléments.
	Ou : il existe 20 façon de choisir 3 éléments parmi 6 en ignorant l'ordre des éléments choisis.
}

\exe{}{
	Une carte bancaire possède un code à 4 chiffres, de 0 à 9.
	\begin{enumerate}
		\item
		Combien y a-t-il
		\begin{enumerate}
			\item
			de codes possibles ?
			\item
			de codes se terminant par un chiffre pair ?
			\item
			de codes contenant au moins une fois 7 ?
			\item
			de codes contenant exactement une fois 7 ?
		\end{enumerate}
		\item
		Une banque force son utilisateur a utiliser 4 chiffres différents pour augmenter la sécurité.
		Combien y a-t-il
		\begin{enumerate}
			\item
			de codes possibles ?
			\item
			de codes se terminant par un chiffre pair ?
			\item
			de codes contenant exactement une fois 7 ?
		\end{enumerate}
		\item
		La banque a-t-elle eu raison de forcer l'utilisation de 4 chiffres différents ?	
	\end{enumerate}
}{exe:bank}{
	\begin{enumerate}
		\item
		Combien y a-t-il
		\begin{enumerate}
			\item
			Un code est une liste $( \_ ; \_ ; \_ ; \_)$ de quatre chiffres de 0 à 9.
			Chaque chiffre a donc 10 possibilités, et au total il y a $10\times10\times10\times10 = 10~000$ combinaisons possibles.
			\item
			Le quatrième chiffre n'a plus que 5 possibilités, ce qui donne $10\times10\times10\times5 = 5~000$ combinaisons possibles.
			\item
			Comptons plutôt les codes ne contenant jamais 7 : il y a en $9^4 = 6~561$ car chaque numéro du code a $9$ possibilités.
			Les codes contenant au moins une fois 7 sont donc au nombre de $10^4 - 9^4 = 3~439$.
			\item
			Comptons les codes ayant un 7 en première positions et aucun autre 7 : $(7 ; \_ ; \_ ; \_)$ avec $9\times9\times9 = 729$ possibilités.
			Les codes ayant un 7 en deuxième position sont au même nombre, idem pour la troisième et quatrième position.
			En ajoutant les résultats nous ne comptons aucun code en double car chaque code contient un unique 7 en des positions différentes.
			Il y a donc $4\times729 = 2~916$ différents codes.
		\end{enumerate}
		\item
		Une banque force son utilisateur a utiliser 4 chiffres différents pour augmenter la sécurité.
		Combien y a-t-il
		\begin{enumerate}
			\item
			On compte ici les arrangements : dans la liste $( \_ ; \_ ; \_ ; \_)$, le premier chiffre a 10 possibilités, le deuxième 9, le troisième 8, et le quatrième 7.
			Il y a donc $10\times9\times8\times7 = 5~040$ codes différents.
			\item
			Au lieu de compter à partir du premier chiffre du code, ce qui changerait potentiellement le nombre de possibilités du dernier (si un pair est choisi), comptons à partir de la fin.
			Le dernier chiffre a 5 possibilités (0 ; 2 ; 4 ; 6 ; 8), l'avant-dernier en a 9 car il est différent du dernier, puis 8, puis 7.
			Donc le nombre de tels codes est $5\times9\times8\times7 = 2~520$.
			\item
			Comptons les codes ayant un 7 en première positions et aucun autre 7 : $(7 ; \_ ; \_ ; \_)$ avec $9\times8\times7 = 504$ possibilités.
			Les codes ayant un 7 en deuxième position sont au même nombre, idem pour la troisième et quatrième position.
			En ajoutant les résultats nous ne comptons aucun code en double car chaque code contient un unique 7 en des positions différentes.
			Il y a donc $4\times507 = 2~016$ différents codes.
		\end{enumerate}
		\item
		D'une part, une attaque « force brute » (c'est-à-dire où l'on teste toutes les combinaisons possibles) est plus efficace plus il y a de combinaisons possibles.
		Pour cela, interdire les chiffres doubles est contre-productif car le nombre de code a été pratiquement divisé par deux : avant répétition, il y a $10^4 = 10~000$ codes possibles, alors que sans répétition il y a $10\times9\times8\times7 = 5~040$ codes possibles.
		
		D'autre part, un code qui contient plusieurs fois le même code (par exemple $(1;1;1;1)$) est plus facile à reconnaître en regardant le propriétaire de la carte le taper.
		Pour cela, interdire les chiffres doubles peut s'avérer utile !
	\end{enumerate}
}

\exe{}{
	On pose un pion sur la case inférieure gauche d'un échiquier $4\times4$.
	On déplace le pion de case en case afin d'atteindre la case supérieure droite.
	Le pion peut soit se déplacer d'une case vers le haut, soit d'une case vers la droite. Il est donc impossible de faire une boucle.
	En combien de façons différentes peut-il atteindre la case supérieure droite ?
}{exe:binom2}{
	Les seuls mouvements du pions sont « Haut », qu'on notera $H$, et « Droite », noté $D$.
	Du point de départ jusqu'au point d'arrivée, le pion doit faire exactement trois fois $H$, et trois fois $D$.
	L'ordre des$H$ et $D$ détermine le chemin pris par le pion : deux mots différents correspondent à deux chemins différents et vice versa.
	Un exemple de mot est « HHDHDD » ; un autre est « DHDHDH ».
	
	Pour compter le nombre de mots différents, il suffit de placer les 3 $H$ dans les 6 positions possibles, les $D$ remplissant les trous.
	Un mot est donc associé à un unique choix de 3 emplacements parmi les 6 possibles dans lesquels mettre un $H$.
	Il y a donc ${6 \choose 3} = 20$ mots et donc chemins possibles. 
}



\exe{difficulty=1}{
	Soient $m \leq n$ deux entiers naturels non nuls.
	Montrer que le nombre de fonctions $f : [m] \rightarrow [n]$ strictement croissantes est égal à ${n \choose m}$.
	
	\emph{Indication : montrer qu'un $m$-arrangement de $[n]$ correspond à une unique fonction $f$ strictement croissante.}
}{exe:strict-increasing}{
	Un choix de $m$ éléments parmi $n$ aboutit à une unique fonction strictement croissante car il n'existe qu'une seule façon de ranger $m$ entiers distincts par ordre croissant.
	
	Plus rigoureusement, l'application qui associe à chaque fonction $f : [m] \rightarrow [n]$ l'ensemble $\{ f(1) ; f(2) ; \dots ; f(m) \} \subseteq [n]$ est une bijection des fonctions strictement croissantes vers les sous-ensembles de $[n]$ contenant $m$ éléments.
	La surjectivité découle de la possibilité de ranger tout ensemble par ordre croissant et d'ainsi définir une fonction pré-image strictement croissante.
	L'injectivité découle de l'unicitié de cet ordre croissant.
}

\exe{difficulty=1}{
	Soit $n\in\N$. En comparant les termes en $x^n$ de l'identité
		$(1+x)^{2n} = (1+x)^n \cdot (1+x)^n, $
	montrer que
		\[ \sum_{k=0}^{n} {n \choose k} {n \choose n-k} = {2n \choose n}. \]
}{exe:id-newton}{
	D'après le binôme de Newton,
		\[ (1+x)^{2n} = \sum_{k=0}^{2n} {2n \choose k} x^k. \]
	Son coeffiicent en $x^n$ est donc bien $2n \choose n$.
	
	De plus,
		\[ (1+x)^n (1+x)^n = \left[ {n \choose 0} + {n \choose 1}x + {n \choose 2}x^2 + \dots + {n \choose n}x^n \right]\left[ {n \choose 0} + {n \choose 1}x + {n \choose 2}x^2 + \dots + {n \choose n}x^n \right]. \]
	Les produits contribuant au terme $x^n$ sont donc
		\[ {n \choose 0} \times {n \choose n}x^n \]
	et 
		\[ {n \choose 1}x \times {n\choose n-1} x^{n-1} \]
	et
		\[ {n \choose 2}x^2 \times {n\choose n-2}x^{n-2} \]
	etc...
	Par conséquent nous trouvons
		\[ {2n \choose n} = {n \choose 0}{n \choose n} + {n \choose 1}{n \choose n-1} + {n \choose 2}{n \choose n-2} + \dots + {n \choose n}{n \choose 0} =  \sum_{k=0}^{n} {n \choose k} {n \choose n-k}. \]
}

\exe{}{
	Calculer le nombre de multiples de $7$, $11$, ou $13$ dans $[1000] = \{ 1 ; 2 ; 3 ; \dots ; 1000 \}$.
	
	Écrire un programme informatique permettant de vérifier son résultat.
}{exe:divisors}{
	Le principe d'inclusion-exclusion permet de trouver 280.

	\input{../../exe:divisors.py}
}

\exe{difficulty=1}{
	Dans l'ensemble $[1000] = \{ 1 ; 2 ; 3 ; \dots ; 1000 \}$, les multiples de 8, de 12, et de 15 sont supprimés.
	Combien de nombres restent-il ?
	
	Écrire un programme informatique permettant de vérifier son résultat.
}{exe:divisors-hard}{
	Le principe d'inclusion-exclusion permet de trouver 783.
	
	\input{../../exe:divisors-hard.py}
}



\exe{difficulty=1}{
	Un ballon de football est constituté de 32 pentagones et hexagones réguliers comme ci-dessus :
	chaque pentagone est entouré de cinq hexagones et chaque hexagone est entouré de trois pentagones et trois hexagones qui s'alternent.

	Pour des pentagones $p$ et hexagones $h$, considérons l'ensemble 
		\[ S = \bigset{ (p, h) \tq \text{ $p$ et $h$ partagent un côté } }. \]
	En comptant $|S|$ de deux façons différentes, montrer que $5p = 3h$ et en déduire le nombre de pentagones et d'hexagones sur le ballon.
}{exe:double-counting}{
	Calculons $|S|$ en parcourant tous les pentagones $p$ et en comptant les couples $(p, h)$ appartenant à $S$.
	Comme chaque pentagone est entouré de cinq hexagones, il existe $5$ couples $(p, h)$ exactement.
	Par conséquent, $|S| = 5p$ car nous ajoutons $p$ fois $5$.
	
	Calculons $|S|$ en parcourant tous les hexagones $h$ et en comptant les couples $(p, h)$ appartenant à $S$.
	Comme chaque hexagone est entouré de trois pentagones, il existe $3$ couples $(p, h)$ exactement.
	Par conséquent, $|S| = 3h$ car nous ajoutons $h$ fois $3$.
	
	Finalement, $5p = 3h$.
	Comme $p+h = 32$ d'après l'énoncé, nous obtenons un système
		\[ \systeme{p+h = 32, \\ 5p = 3h.} \]
	dont nous déduisons que $3p + 3h = 96$ et donc $3p + 5h = 96$ et $h = \frac{96}{8} = 12$.
	Puis $h = 32-p = 20$.
}

\exe{}{
	 Une secrétaire dispose de $12$ enveloppes portant chacune l'adresse d'un client différent. Elle a aussi $12$ lettres différentes à envoyer à chacun des $12$ clients. Ne sachant pas pour chaque client quelle lettre doit lui être destinée, elle met au hasard chacune des $12$ lettres dans chacune des $12$ enveloppes et les poste en priant le ciel pour que chaque client reçoive sa lettre.
	 \begin{enumerate}
	  	\item 
	  	Quelle est la probabilité que son vœu soit exaucé ?
	  	\item 
	  	Quelle est la probabilité qu'aucun des clients ne reçoivent sa lettre ?
%	  	\item 
%	  	Quelle est la probabilité que deux clients exactement reçoivent leur lettre ?
	 \end{enumerate}
}{exe:derangement}{
	 \begin{enumerate}
		  \item 
		  Il y a une seule permutation favorable pour que tous les clients reçoivent leur lettre et il y a $12!=479001600$ permutations possibles. La probabilités que tous les clients reçoivent leurs lettres est donc : $\dfrac{1}{12!}=\dfrac{1}{479001600}.$
	  	\item 
	  	La probabilité qu'aucun client ne reçoive sa lettre est égal au nombre de dérangements d'un ensemble à $12$ éléments divisé par $12!$ : 
			\[  \sum_{k=0}^{12}\frac{(-1)^k}{k!}=\frac{16019531}{43545600}\approx 0,368. \]
%	  	\item 
%	  	La probabilité que deux clients exactement reçoivent leur lettre est la probabilité que la secrétaire n'ait pas permuté leurs lettres en les mettant dans l'enveloppe tout en ayant dérangé les autres, c'est-à-dire :  
%			\[ \frac{1}{2!}\times \frac{10!\sum_{k=0}^{10}\frac{(-1)^k}{k!}}{10!}\approx 0,184 \ (\approx \frac{1}{2!e}). \]
	 \end{enumerate}
}

\exe{}{
	Donner un ensemble de 6 entiers relatifs tels qu'aucune différence de deux d'entre eux n'est multiple de 6.

	Montrer que parmi 7 entiers relatifs, la différence de deux d'entre eux est un multiple de 6.
}{exe:pidgin-whole0}{
	L'ensemble $\{ 0 ; 1 ; 2 ; 3 ; 4 ; 5 \}$ fonctionne : aucune différence n'est un multiple de 6.
	
	
	Pour raisonner plus généralement sur sept entiers, nous utilisons le principe des chaussettes dans les tiroirs : en mettant $n+1$ chaussettes dans $n$ tiroirs, un tiroir contient nécessairement plus de deux chaussettes.
	
	Comme il existe six restes possibles modulo 6, parmi les sept entiers relatifs, deux d'entre eux auront nécessairement le même reste : ils sont congrus modulo 6.
	Leur différence est donc congrue à 0 modulo 6, c'est-à-dire que c'est un multiple de 6.
}

\exe{}{
	Montrer que parmi 100 personnes, 9 d'entre elles sont nécessairement nées le même mois de l'année.
}{exe:pidgin-whole1}{
	Nous utilisons le principe des chaussettes dans les tiroirs généralisé.
	
	Rangeons chaque personne dans une boite correspondant à son mois de naissance.
	Supposons qu'aucun mois ne contienne 9 personnes : chaque mois contient donc moins de 8 personnes.
	Au total, il y aura donc moins de $12\times 8 = 96$ personnes rangées.
	Ceci n'est pas possible car 100 personnes ont été rangées.
}


\exe{difficulty=1}{
	Montrer qu'il existe un nombre de la forme $N = 11...100...0 \in\N$ (que des $1$ puis que des $0$) qui est multiple de $23$ et inférieur ou égal à $10^{24}$.
	
	\emph{Exemple : 111100 qui n'est malheureusement pas multiple de 23.}
}{exe:pidgin-whole2}{
	Considérons la suite d'entiers
		\[ \bigset{ 0 ; 1 ; 11 ; 111 ; 1111 ; \dots ; \underbrace{11...11}_{\text{$23$ fois}} }, \]
	et leur reste modulo 23.
	Il existe 23 restes différents, mais 24 nombres dans l'ensemble.
	Le principe des chaussettes dans les tiroirs implique que deux éléments partagent le même reste modulo 23.
	Leur différence (grand moins petit) est donc un multiple de 23 et a la forme requise.
	La plus grande telle différence est $\underbrace{11...11}_{\text{$23$ fois}} - 0 \leq 10^{24}$ comme requis.
}

\exe{difficulty=1}{
	Soit $n\in\N^*$ et $n = k_1 + k_2 + \dots + k_m$ une partition de $n$ en $m$ entiers naturels.
	
	Montrer que $k_1! k_2! \cdots k_m!$ divise $n$.
	
	Le nombre 
		\[ \frac{n!}{k_1! k_2! \cdots k_m!} = {n \choose k_1, k_2, \dots, k_m} \in \N \]
	est un \emph{coefficient multinomial}.
}{exe:multinomial}{
	Prenons les éléments de $[n] = \{ 1 ; 2 ; \dots ; n \}$ et découpons-les en $m$ ensembles de taille $k_1, k_2, \dots, k_m$.
	Comptons le nombre d'ensembles distincts obtenus.
	Pour cela, partons de listes ordonnées, et divisons par le nombre de permutation à l'interieur de chaque liste pour compter le nombre d'ensembles.
	
	Prenons une permutations de $[n]$. Il en existe $n!$.
	Les $k_1$ premiers éléments sont comptés un facteur $k_1!$ en trop si l'on ne souhaite pas garder l'ordre.
	Les $k_2$ suivants sont comptés un facteur $k_2!$ en trop.
	Etc...
	Au total, le nombre de partitionnements de $[n]$ en sous-ensembles (donc sans regarder l'ordre) est égal à 
		\[  \frac{n!}{k_1! k_2! \cdots k_m!} = {n \choose k_1, k_2, \dots, k_m} \in \N. \]
		
	Remarque : le coefficient multinomial apparaît lorsqu'on souhaite développer $(1+x+y)^n$ puis $(1+x+y+z)^n$, etc...
}



%%%%%%%%%%%%%%%%%%%%%%%%
%%%%%%%%%% SET 2  %%%%%%%%%%
%%%%%%%%%%%%%%%%%%%%%%%% 


\exe{}{
	Nous souhaitons mettre $n$ boules identiques dans $k$ urnes.
	Pour cela, nous représentons les boules avec des étoiles $\star$ et les séparations entre urnes avec des barres $|$.
	
	Chaque façon de ranger les boules correspond à un unique mot sur l'alphabet $\bigset{ \star ; | }$
	Par exemple, le mot « $\star|\star||\star\star$ » correspond à ranger une boule dans la première urne, une boule dans la deuxième, zéro dans la troisième, et deux dans la dernière.
	\begin{enumerate}
		\item
		Combien d'étoiles y a-t-il ?
		\item
		Combien de barres y a-t-il ?
		\item
		Combien de façons différentes existe-t-il de mettre $n$ boules identiques dans $k$ urnes ?
	\end{enumerate}
}{exe:stars-bars-bins}{
	
}


\exe{}{
	Considérons la fonctions indicatrice d'Euler
		\[ \phi(n) = \left|\bigset{ 1 \leq k \leq n \tq \pgcd(k, n) = 1 } \right|. \]
	Montrer grâce au principe d'inclusion-exclusion que
		\[ \phi(n) = n \prod_{p|n} \left(1-\frac1p\right), \]
	où le produit est pris sur les nombres premiers $p$ divisant $n$.
}{exe:Euler-totient}{

}

\exe{}{
	Quatre personnes ont chacune un cadeau à offrir au hasard.
	Comme elles ne peuvent pas toutes se retrouver au même endroit, elles se divisent en deux groupes (non vides) et s'offrent les cadeaux au sein de chaque groupe. 
	Pour ne pas se retrouver avec son propre cadeau, la première personne offre son cadeau à la deuxième, qui offre son cadeau à la troisième, etc... jusqu'à la dernière personne du groupe qui offre son cadeau à la première.
	
	Combien y a-t-il de permutations différentes des cadeaux ?
	Les lister toutes.
}{exe:Stirling1}{

}

\exe{}{
	Rappeler les valeurs des nombres de Stirling de première espèce $|s(n,1)|$ et $|s(n,n)|$ pour $n\geq1$.

	À l'aide de la relation de récurrence
		\[ |s(n+1, k)| = |s(n,k-1)| + n|s(n,k)|, \]
	remplir le triangle ci-dessous avec les valeurs de $|s(n,k)|$.

	\begin{center}
	\scalebox{1.5}{
	\begin{tabular}{>{$}c<{$}|*{7}{c}}
		\multicolumn{1}{l}{$n$} &&&&&\\\cline{1-1} 
		1 \\\hline
		2 \\\hline
		3 \\\hline
		4 \\\hline
		5 \\\hline
		6 \\\hline
		7 \\\hline
		\multicolumn{1}{l}{} &1&2&3&4&5&6&7\\\cline{2-7}
		\multicolumn{1}{l}{} &\multicolumn{7}{c}{$k$}
	\end{tabular}
	}
	\end{center}
}{exe:Stirling-1ere}{
	
}

\exe{}{
	Développer et réduire le produit
		\[ p(x) = x(x+1)(x+2)(x+3)(x+4)(x+5)(x+6) \]
	à l'aide de l'exercice \ref{exe:Stirling-1ere}.
	
	En déduire le développement de  $q(x) = x(x-1)(x-2)(x-3)(x-4)(x-5)(x-6)$.
	
	\emph{Indication : remplacer $x$ par $-x$.}
}{exe:Stirling-1ere-gen}{

}


\exe{}{
	Six personnes ont chacune un cadeau à offrir au hasard.
	Comme elles ne peuvent pas toutes se retrouver au même endroit, elles se divisent en deux groupes (non vides) et s'offrent les cadeaux au sein de chaque groupe. 
	Pour ne pas se retrouver avec son propre cadeau, la première personne offre son cadeau à la deuxième, qui offre son cadeau à la troisième, etc... jusqu'à la dernière personne du groupe qui offre son cadeau à la première.
	
	Combien y a-t-il de permutations différentes des cadeaux ?
}{exe:Stirling2}{

}

\exe{}{
	Rappeler les valeurs des nombres de Stirling de seconde espèce $S(n,1)$ et $S(n,n)$ pour $n\geq1$.

	À l'aide de la relation de récurrence
		\[ S(n+1, k) = S(n,k-1) + k S(n,k), \]
	remplir le triangle ci-dessous avec les valeurs de $S(n,k)$.

	\begin{center}
	\scalebox{1.5}{
	\begin{tabular}{>{$}c<{$}|*{7}{c}}
		\multicolumn{1}{l}{$n$} &&&&&\\\cline{1-1} 
		1 \\\hline
		2 \\\hline
		3 \\\hline
		4 \\\hline
		5 \\\hline
		6 \\\hline
		7 \\\hline
		\multicolumn{1}{l}{} &1&2&3&4&5&6&7\\\cline{2-7}
		\multicolumn{1}{l}{} &\multicolumn{7}{c}{$k$}
	\end{tabular}
	}
	\end{center}
}{exe:Stirling-2nde}{
	
}




